% ============================================================================
% Մանկավարժական լուծումներ - Դիսկրետ մաթեմատիկա, Գործնական 6
% (Զուգորդություններ և բազմանդամային գործակիցներ). Խնդիրներ 8, 9, 11.
% Հայերեն տարբերակ, XeLaTeX (fontspec + Sylfaen). Կազմել՝ xelatex-ով երկու անգամ։
% Տերմինաբանությունը՝ glossary_discrete_focused.md / glossary_Apr.csv-ի համաձայն։
% ============================================================================
\documentclass[11pt,a4paper]{article}

% ---- Տառատեսակ (XeLaTeX) ----
\usepackage{fontspec}
\setmainfont{Sylfaen}

% Հայերենը plain XeLaTeX-ում տողադարձման կանոններ չունի, ուստի երկար տողերը չեն
% կարող կոտրվել և դուրս կգան լուսանցքից։ Թող TeX-ը ձգի բառամիջյան բացատները։
\tolerance=2000
\emergencystretch=3.5em

% ---- Մաթեմատիկա ----
\usepackage{amsmath,amssymb,amsthm}
\usepackage{mathtools}

% ---- Դասավորություն ----
\usepackage[margin=2.3cm]{geometry}
\usepackage{enumitem}
\usepackage{booktabs}
\usepackage{array}

% ---- Գրաֆիկա ----
\usepackage{xcolor}
\usepackage[most]{tcolorbox}
\usepackage{tikz}
\usetikzlibrary{arrows.meta,positioning,calc,fit,backgrounds,shapes.misc}
\usepackage{pgfplots}
\pgfplotsset{compat=1.18}
\usepackage{hyperref}
\hypersetup{colorlinks=true,linkcolor=blue!55!black,urlcolor=blue!55!black}

% ---- Գունապնակ (Հայաստանի դրոշի գույները) ----
\definecolor{armRed}{HTML}{D90012}
\definecolor{armBlue}{HTML}{0033A0}
\definecolor{armOrange}{HTML}{F2A800}
\definecolor{ink}{HTML}{22303C}

\setlength{\parindent}{0pt}
\setlength{\parskip}{5pt}

% ---- Վանդակներ ----
\newtcolorbox{problembox}[1][]{enhanced, breakable, sharp corners=south,
  colback=armOrange!10, colframe=armOrange!75!black, boxrule=0.8pt, arc=2mm,
  fonttitle=\bfseries\color{white}, coltitle=white,
  attach boxed title to top left={xshift=4mm,yshift=-2.4mm},
  boxed title style={colback=armOrange!85!black, arc=1mm},
  title={Խնդիր #1}}

\newtcolorbox{ideabox}{enhanced, breakable,
  colback=armBlue!6, colframe=armBlue!70!black, boxrule=0.8pt, arc=2mm,
  fonttitle=\bfseries\color{white}, coltitle=white,
  attach boxed title to top left={xshift=4mm,yshift=-2.4mm},
  boxed title style={colback=armBlue!75!black, arc=1mm},
  title={Գաղափարը}}

\newtcolorbox{methodbox}{enhanced, breakable,
  colback=gray!8, colframe=gray!55!black, boxrule=0.7pt, arc=2mm,
  fonttitle=\bfseries, title={$\bigstar$\ Հիշարժան մեթոդ}}

\newtcolorbox{answerbox}{enhanced, sharp corners=south,
  colback=green!8, colframe=green!50!black, boxrule=1pt, arc=2mm, left=4mm,right=4mm}
\newcommand{\answer}[1]{\begin{answerbox}\centering\textbf{Պատասխան՝}\quad #1\end{answerbox}}

% ---- Վերնագիր ----
\title{\vspace{-1.2cm}\textbf{Դիսկրետ մաթեմատիկա - Գործնական 6}\\[2pt]
\large Զուգորդություններ, բինոմիալ և բազմանդամային գործակիցներ՝ մանրամասն լուծումներ\\[2pt]
\normalsize\itshape դատողությունները՝ քայլ առ քայլ բացված}
\author{}
\date{}

\begin{document}
\maketitle
\vspace{-1.0cm}

\section*{Նախապատրաստում}
Այս թերթիկը \emph{ընտրությունների և դասավորությունների հաշվման} մասին է։ Ամբողջ գլուխը հենվում է
չորս բազմակիրառ գաղափարի վրա՝ \textbf{արտադրյալի կանոն} (անկախ ընտրությունները բազմապատկվում են),
\textbf{զուգորդություններ} $\binom{n}{r}$ ($n$ տարրից չկարգավորված $r$-ենթաբազմություն ընտրելը),
\textbf{կրկնություններով զուգորդություններ} (նույնական տարրեր բաշխել պիտակավորված արկղերի մեջ, այլ
կերպ՝ աստղեր և ձողիկներ) և \textbf{բազմանդամային գործակից} $\frac{n!}{n_1!\,n_2!\cdots}$ (դասավորել
մուլտիբազմություն, որտեղ որոշ տարրեր կրկնվում են)։ Արվեստն այն է, որ ճանաչես, թե որ մոդելն է
իրականում տվյալ խնդիրը։ Ստորև լուծում ենք 8, 9 և 11 խնդիրները, և յուրաքանչյուրի համար նախ բացատրում
ենք, թե \emph{ինչու} է ընտրված մոդելը համապատասխանում, ապա նոր պտտում ենք բռնակը։

% ============================================================================
\section*{Խնդիր 8 - Ցանցային ճանապարհներ ստուգակետով}

\begin{problembox}[8]
Ամբողջաթիվ ցանցում հաշվիր $O(0,0)$-ից $A(7,8)$ տանող միապաղաղ ճանապարհները, որոնք անցնում են $(3,4)$
կետով։ Յուրաքանչյուր քայլ գնում է մեկ միավոր \textbf{աջ} կամ մեկ միավոր \textbf{վեր}։
\end{problembox}

\begin{ideabox}
Նախ՝ ինչու՞ են միապաղաղ ճանապարհներն ընդհանրապես հաշվվում բինոմիալ գործակցով։ $O(0,0)$-ից $A(7,8)$
հասնելու համար՝ օգտագործելով միայն աջ ($R$) և վեր ($U$) քայլեր, պետք է անես ճիշտ $7$ աջ և $8$ վեր
քայլ՝ \emph{ինչ-որ կարգով}։ Այսինքն՝ ճանապարհը \textbf{հենց} $7$ $R$-ից և $8$ $U$-ից կազմված բառ է,
$7+8=15$ քայլերի հաջորդականություն։ Մեկը կառուցելու համար պարզապես պետք է որոշես, թե $15$ դիրքերից որ
$7$-ն են պահում $R$-երը (մնացած դիրքերը հարկադրաբար $U$-եր են)՝ դա $\binom{15}{7}$ տարբերակ է։
Ճանապարհներ հաշվելը դառնում է հաշվել, թե որտեղ են գնում աջերը։

Հիմա ստուգակետը։ $(3,4)$-ով անցնող ցանկացած ճանապարհ մաքուր կերպով բաժանվում է երկու կեսի՝
$O\to(3,4)$ ճանապարհորդությունը, ապա $(3,4)\to A$ ճանապարհորդությունը։ Այս երկու կեսերն \textbf{անկախ}
են (ինչ էլ որ անես $(3,4)$-ից առաջ՝ չի սահմանափակում, թե ինչ կանես հետո), ուստի արտադրյալի կանոնով
առաջին կեսերի քանակը բազմապատկում ենք երկրորդ կեսերի քանակով։ Ստուգակետը մեկ դժվար հաշիվը վերածում է
երկու հեշտի։
\end{ideabox}

\textbf{Քայլ 1 - Առաջին հատված՝ $O(0,0)\to(3,4)$։}
Այս հատվածը պահանջում է $3$ աջ և $4$ վեր քայլ, ուստի՝ $3+4=7$ քայլ։ Ճանապարհն այն ընտրությունն է, թե
այդ $7$ քայլերից որ $3$-ն են աջերը՝
$$\binom{7}{3}=\frac{7!}{3!\,4!}=35.$$

\textbf{Քայլ 2 - Երկրորդ հատված՝ $(3,4)\to A(7,8)$։}
Այս հատվածը պահանջում է $7-3=4$ աջ և $8-4=4$ վեր քայլ, ուստի՝ $8$ քայլ։ Ընտրիր, թե որ $4$-ն են աջերը՝
$$\binom{8}{4}=\frac{8!}{4!\,4!}=70.$$

\textbf{Քայլ 3 - Բազմապատկենք (արտադրյալի կանոն)։}
$$\binom{7}{3}\binom{8}{4}=35\cdot 70=2450.$$

\textbf{Ստուգում։} Ողջամիտ սահման՝ $O\to A$ ճանապարհների \emph{ընդհանուր} քանակը (անտեսելով
ստուգակետը) $\binom{15}{7}=6435$ է, և մեր պատասխանը՝ $2450$-ը, հարմարավետ կերպով ավելի փոքր է,
ինչպես և պետք է, քանի որ ոչ բոլոր ճանապարհներն են այցելում $(3,4)$։ Նաև բաժանումը համաձայնեցված է՝
$35\cdot 70 = 2450$, և բոլոր $6435$ ճանապարհների ուղիղ համակարգչային թվարկումը գտնում է ճիշտ
$2450$-ը, որ անցնում են $(3,4)$-ով։

\answer{$\displaystyle \binom{7}{3}\binom{8}{4}=35\cdot 70=2450$}

\begin{methodbox}
\textbf{Փոխանցելի տեխնիկա - ճանապարհներ և ստուգակետով բաժանումը՝} $(0,0)$-ից $(a,b)$ տանող միապաղաղ
ցանցային ճանապարհը համապատասխանում է $a$ աջից և $b$ վերից կազմված բառի, ուստի դրանք $\binom{a+b}{a}$
հատ են։ Ճանապարհը ֆիքսված $P$ կետով անցկացնելու համար կտրիր այն $P$-ում և \textbf{բազմապատկիր} երկու
անկախ հատվածների քանակները։ Կետն \textbf{արգելելու} համար հանիր՝ ($P$-ով անցնող ճանապարհները)
ընդհանուրից։
\end{methodbox}

\medskip
\textbf{Բաժանումը տեսնելը։} Ստորև ցանցը ցույց է տալիս երկու հատվածները։ Յուրաքանչյուր հատված իր
սեփական փոքրիկ ճանապարհ-հաշվելու խնդիրն է՝ $(3,4)$ ստուգակետը զոդում է դրանք իրար։

\begin{center}
\begin{tikzpicture}[scale=0.62, font=\small]
  \draw[step=1, gray!35, very thin] (0,0) grid (7,8);
  \draw[->] (0,0) -- (7.6,0) node[right]{$x$};
  \draw[->] (0,0) -- (0,8.6) node[above]{$y$};
  \fill[armBlue] (0,0) circle (4pt) node[below left]{$O(0,0)$};
  \fill[armOrange!90!black] (3,4) circle (4pt) node[above left]{$(3,4)$};
  \fill[armRed] (7,8) circle (4pt) node[above right]{$A(7,8)$};
  \draw[armBlue, very thick, line cap=round]
     (0,0) -- (1,0) -- (1,1) -- (2,1) -- (2,2) -- (2,3) -- (3,3) -- (3,4);
  \draw[armRed, very thick, line cap=round]
     (3,4) -- (4,4) -- (4,5) -- (5,5) -- (6,5) -- (6,6) -- (6,7) -- (7,7) -- (7,8);
  \node[armBlue]  at (1.4,5.0) {հատված 1՝ $\binom{7}{3}=35$};
  \node[armRed]   at (5.6,3.1) {հատված 2՝ $\binom{8}{4}=70$};
\end{tikzpicture}
\end{center}

% ============================================================================
\section*{Խնդիր 9 - 6-ի տրոհումը երեք գումարելիի}

\begin{problembox}[9]
Հաշվիր $6$-ը երեք գումարելիի գումարի տեսքով գրելու տարբերակների քանակը։
\end{problembox}

\begin{ideabox}
\emph{Տրոհում} բառն ինքնին կարող է նշանակել \emph{չկարգավորված} բաժանում, ուստի քանակն իսկապես կախված
է պայմանավորվածությունից։ Այս վարժությունը գտնվում է \textbf{կրկնություններով զուգորդությունների}
բաժնում, ինչը ֆիքսում է ընթերցումը՝ երեք գումարելիները \textbf{կարգավորված} են (առաջին, երկրորդ և
երրորդ տեղերը տարբեր են) և \textbf{զրոն թույլատրելի է}, ճիշտ ինչպես դասախոսության օրինակում, որտեղ
$4=x_1+x_2$-ը հաշվվում էր $5$ տարբերակ՝ $4{+}0,\,3{+}1,\,2{+}2,\,1{+}3,\,0{+}4$։ Այսպիսով՝ հաշվում ենք
$$x_1+x_2+x_3=6,\qquad x_i\ge 0$$
հավասարման ամբողջ լուծումները։ Ինչու՞ է սա \textbf{աստղեր և ձողիկներ} խնդիր։ Պատկերացրու $6$-ը որպես
շարքում վեց նույնական աստղ։ Երեք գումարելիների ընտրությունը նույնն է, ինչ աստղերի մեջ \textbf{երկու
ձողիկ} գցելը՝ շարքը երեք խմբի կտրելու համար՝ առաջին ձողիկից առաջ աստղերը $x_1$ են, ձողիկների միջև՝
$x_2$, երկրորդ ձողիկից հետո՝ $x_3$։ Դատարկ խմբեր թույլատրվում են (երկու ձողիկ կարող են կից լինել՝
տալով զրո գումարելի)։ Երկու ձողիկների յուրաքանչյուր տեղադրություն տալիս է ճիշտ մեկ լուծում, և
հակառակը, ուստի պարզապես հաշվում ենք ձողիկների տեղադրությունները։
\end{ideabox}

\textbf{Քայլ 1 - Կոդավորենք որպես սիմվոլների շարք։}
Շարքում դասավորիր $6$ աստղ և $2$ ձողիկ, օրինակ՝ $\star\star\,|\,\star\,|\,\star\star\star$, որը
կարդացվում է $x_1=2,\ x_2=1,\ x_3=3$։ Ամբողջ շարքն ունի $6+2=8$ սիմվոլային դիրք։

\textbf{Քայլ 2 - Հաշվենք դասավորությունները։}
Լուծումը ֆիքսվում է, հենց որոշում ենք, թե $8$ դիրքերից որ $2$-ն են ձողիկները (մնացածն աստղեր են)՝
$$\binom{6+3-1}{3-1}=\binom{8}{2}=\frac{8\cdot 7}{2}=28.$$

\textbf{Ստուգում։} $x_1+x_2+x_3=6$-ի $x_i\ge0$ լուծումների ուղիղ հաշիվ՝ երբ $x_1$-ը փոխվում է $0$-ից
$6$, մնացած $x_2+x_3=6-x_1$-ն ունի $(6-x_1)+1$ լուծում, ինչը տալիս է $7+6+5+4+3+2+1=28$։ Ուղիղ եռակի
ցիկլը նույնպես վերադարձնում է $28$։ \checkmark

\answer{$\displaystyle \binom{8}{2}=28$}

\begin{methodbox}
\textbf{Փոխանցելի տեխնիկա - աստղեր և ձողիկներ՝} $x_1+\dots+x_k=n$-ի կարգավորված ոչ բացասական ամբողջ
լուծումների քանակը $\binom{n+k-1}{k-1}$ է ($n$ աստղերի մեջ տեղադրիր $k-1$ ձողիկ)։ Սա
\emph{կրկնություններով զուգորդություն} է՝ այն հավասար է $k$ տեսակից $n$ չափի մուլտիբազմություն ընտրելու
տարբերակների քանակին։ Եթե փոխարենը յուրաքանչյուր $x_i\ge 1$, ապա նախ յուրաքանչյուր փոփոխականին տուր
մեկ միավոր, ապա լուծիր $x_1'+\dots+x_k'=n-k$, ինչը տալիս է $\binom{n-1}{k-1}$։
\end{methodbox}

\medskip
\textbf{Ձողիկների պատկերը։} Վեց աստղի մեջ երկու ձողիկի երեք օրինակ տեղադրություն՝

\begin{center}
\begin{tikzpicture}[font=\small, star/.style={armRed, scale=1.2}, bar/.style={armBlue, very thick}]
  \node at (-1.6,0) {$(2,1,3)$՝};
  \foreach \x in {0,1} {\node[star] at (\x*0.5,0) {$\star$};}
  \draw[bar] (1.0,-0.22) -- (1.0,0.22);
  \node[star] at (1.5,0) {$\star$};
  \draw[bar] (2.0,-0.22) -- (2.0,0.22);
  \foreach \x in {0,1,2} {\node[star] at (2.5+\x*0.5,0) {$\star$};}
  \node at (-1.6,-0.8) {$(0,4,2)$՝};
  \draw[bar] (0.0,-0.8-0.22) -- (0.0,-0.8+0.22);
  \foreach \x in {0,1,2,3} {\node[star] at (0.5+\x*0.5,-0.8) {$\star$};}
  \draw[bar] (3.0,-0.8-0.22) -- (3.0,-0.8+0.22);
  \foreach \x in {0,1} {\node[star] at (3.5+\x*0.5,-0.8) {$\star$};}
  \node at (-1.6,-1.6) {$(6,0,0)$՝};
  \foreach \x in {0,1,2,3,4,5} {\node[star] at (\x*0.5,-1.6) {$\star$};}
  \draw[bar] (3.0,-1.6-0.22) -- (3.0,-1.6+0.22);
  \draw[bar] (3.12,-1.6-0.22) -- (3.12,-1.6+0.22);
\end{tikzpicture}
\end{center}
Կից ձողիկները (վերջին շարք) տալիս են զրո գումարելի, ինչը պատճառն է, որ $\binom{8}{2}=28$, ոչ թե ավելի
փոքր $\binom{5}{2}=10$, որը կտար բոլոր-դրական պայմանավորվածությունը՝ իսկ ավելի փոքր է $3$-ը, եթե երեք
գումարելիները դիտվեին \emph{չկարգավորված} ($\{1,1,4\},\{1,2,3\},\{2,2,2\}$)։ Երեք ընթերցում, երեք
թիվ ($28$, $10$, $3$)՝ այստեղ հաշվում ենք կարգավորված եռյակներ՝ զրոներ թույլատրված։

% ============================================================================
\section*{Խնդիր 11 - Բառի տառերի դասավորությունները}

\begin{problembox}[11]
Քանի՞ տարբեր «բառ» (բոլոր տառերի դասավորություն) կարելի է կազմել հայերեն \textbf{մաթեմատիկա} բառի
տառերից։
\end{problembox}

\begin{ideabox}
Եթե բոլոր $10$ տառերը \emph{տարբեր} լինեին, պատասխանը պարզապես $10!$ դասավորություն կլիներ։ Հնարքն այն
է, որ որոշ տառեր \textbf{կրկնվում} են, և նույն տառի երկու օրինակ տեղափոխելը նոր տեսանելի բառ չի տալիս։
Ուստի $10!$-ը ավելորդ է հաշվում՝ այն կրկնվող օրինակները դիտում է որպես տարբերակելի։

Ուղղումն այն է, որ \textbf{բաժանենք} ավելորդ հաշվումը։ Եթե տառը հանդիպում է $m$ անգամ, ապա
յուրաքանչյուր իսկապես տարբեր բառի համար $10!$ հաշիվը թաքնված կերպով այն թվարկում է $m!$ անգամ (մեկ՝ այդ
նույնական օրինակները իրենց դիրքերում տեղափոխելու յուրաքանչյուր եղանակի համար)։ Յուրաքանչյուր կրկնվող
տառի համար $m!$-ի վրա բաժանելը հեռացնում է հենց այդ ավելորդությունը։ Այս հարաբերությունը
\textbf{բազմանդամային գործակիցն} է։

\textbf{Վճռորոշ քայլը մուլտիբազմությունը ճիշտ որոշելն է։} Հայերեն \textbf{մաթեմատիկա} բառում երկու
տարբեր «տ»-հնչյուն կա, որոնք \emph{տարբեր հայերեն տառեր} են՝ շնչեղ \textbf{թ}-ն (երրորդ տառը) և պարզ
\textbf{տ}-ն (յոթերորդ տառը)։ Յուրաքանչյուրը հանդիպում է միայն \emph{մեկ} անգամ, ուստի դրանք
\textbf{կրկնվող} տառ չեն։ Դրանք շփոթելը (ինչպես կաներ լատինատառ \texttt{matematika} գրությունը՝ իր
երկու նույնական \texttt{t}-երով) տալիս է սխալ պատասխան։
\end{ideabox}

\textbf{Քայլ 1 - Թվարկենք տառերի մուլտիբազմությունը։}
Տասը հայերեն տառերը կարգով գրելով և հավասարները խմբավորելով՝
\begin{center}
\small
\begin{tabular}{@{}lcl@{}}
\toprule
տառ & քանակ & նշում\\
\midrule
մ  & $2$ & կրկնվող\\
ա  & $3$ & կրկնվող\\
թ (շնչեղ) & $1$ & առանձին տառ\\
ե  & $1$ &\\
տ (պարզ) & $1$ & \emph{այլ} առանձին տառ\\
ի  & $1$ &\\
կ  & $1$ &\\
\bottomrule
\end{tabular}
\end{center}
Ընդամենը՝ $2+3+1+1+1+1+1=10$ տառ։ Միայն \textbf{մ}-ն (երկու անգամ) և \textbf{ա}-ն (երեք անգամ) են
կրկնվում։

\textbf{Քայլ 2 - Կիրառենք բազմանդամային գործակիցը։}
$$\frac{10!}{2!\,3!\,1!\,1!\,1!\,1!\,1!}=\frac{10!}{2!\,3!}
   =\frac{3\,628\,800}{2\cdot 6}=\frac{3\,628\,800}{12}=302\,400.$$

\textbf{Ստուգում։} Երկու անկախ ստուգում։
\emph{(i) Կառուցենք բառը՝ տառերը $10$ դիրքերի մեջ տեղադրելով։} Ընտրիր \textbf{մ}-երի $2$ դիրքը՝
$\binom{10}{2}=45$։ Մնացած $8$ դիրքից ընտրիր \textbf{ա}-երի $3$-ը՝ $\binom{8}{3}=56$։ Վերջին $5$
տարբեր տառերը լրացնում են մնացած $5$ դիրքը $5!=120$ տարբերակով։ Ապա $45\cdot 56\cdot 120 = 302\,400$՝
համընկնում է։ \emph{(ii)} Մուլտիբազմության տարբեր տեղադրությունների համակարգչային թվարկումը վերադարձնում
է ճիշտ $302\,400$։

\answer{$\displaystyle \frac{10!}{2!\,3!}=302\,400$}

\medskip
\textbf{Ինչու՞ է կարևոր տառերի հաշիվը։} Եթե սխալմամբ օգտագործվեր լատինատառ \texttt{matematika}
գրությունը, երկու \texttt{t}-երը կմիաձուլվեին մեկ կրկնվող տառի, և քանակը կնվազեր մինչև
$$\frac{10!}{2!\,2!\,3!}=151\,200,$$
ինչը ճիշտ կեսն է։ $2$ գործոնը կեղծ նույնական t-երի զույգի $2!$-ն է։ Խնդիրը հայերեն գրության մասին է,
որտեղ երկու t-տառերը տարբեր են, ուստի պատասխանը $302\,400$ է։

\begin{methodbox}
\textbf{Փոխանցելի տեխնիկա - մուլտիբազմության տեղադրություններ՝} $n$ տարր դասավորելու համար, որտեղ
տարբեր տեսակները հանդիպում են $n_1,n_2,\dots,n_r$ պատիկներով (այնպես որ $\sum n_i=n$), տարբեր
դասավորությունների քանակը բազմանդամային գործակիցն է՝
$$\binom{n}{n_1,n_2,\dots,n_r}=\frac{n!}{n_1!\,n_2!\cdots n_r!}.$$
Սկսիր $n!$-ից (բոլոր տարրերը տարբեր), ապա բաժանիր $n_i!$-ի վրա յուրաքանչյուր տեսակի համար՝
չեղարկելու դրա օրինակների չտարբերակվող տեղադրությունները։ Ամենասխալվող քայլը մուլտիբազմությունը ճիշտ
կարդալն է՝ հաշվիր այն, ինչը \emph{իրականում} նույն սիմվոլն է, ոչ թե այն, ինչը տառադարձության մեջ նման
է թվում։
\end{methodbox}

\medskip
\textbf{Նույն բառը հաշվելը երեք մեթոդով (բոլորը համընկնում են)։}

\begin{center}
\small
\begin{tabular}{@{}lll@{}}
\toprule
մեթոդ & հաշվարկ & արժեք\\
\midrule
բազմանդամային        & $10!/(2!\,3!)$                        & $302\,400$\\
դիրք-լրացում         & $\binom{10}{2}\binom{8}{3}\,5!=45\cdot56\cdot120$ & $302\,400$\\
ուղիղ թվարկում       & մուլտիբազմության տարբեր տեղադրությունների հաշիվ  & $302\,400$\\
\bottomrule
\end{tabular}
\end{center}

\end{document}
